
\section{Key findings by study}\label{app:findings}
\begin{description}
\item[Study 1 (pilot).] On 2{,}000-reviewed-positive pilot data, PepCNN
reached AUROC 0.95 within minutes, establishing the pipeline end to end;
the annealing designer produced its first 24 candidates here.
\item[Study 7 (scale-up).] Cluster-aware splitting cut apparent
performance relative to a random split, the first signal that homology
inflates random-split estimates; the frozen screen ensemble dates from
this study and is never retrained afterward.
\item[Study 8 (screen).] 852 reviewed uncharacterized proteins scored;
41 pass the novelty floor; the identical P0ACW4/P0ACW5 sequence pair tops
the list.
\item[Study 9 (APD6).] The frozen ensemble scores APD6-2024 natural AMPs
high on average, a sanity check that the model prefers known AMPs to
uncharacterized proteins, without constituting validation of any
individual candidate.
\item[Study 10 (Veltri base).] On the authors' exact splits, PepCNN
reaches AUC 95.77, the RF 93.32, the GNN 88.18; composition carries much
of the benchmark signal.
\item[Study 11 (stack).] Stacking with tune-partition weights reaches ACC
91.22, MCC 0.8246, AUC 96.61: near-parity with the published AMP Scanner
row, behind ACEP; reported as descriptive, without a paired test.
\item[Study 12 (candidates).] Nine named candidates with novelty and
physicochemical audits and explicit falsifiable MIC predictions.
\item[Study 13 (filtered CV).] A Feb2020 CV that silently filtered records
produced numbers not comparable to the published ones; retained as an
ablation and superseded by study 14.
\item[Study 14 (full CV).] All 4{,}042 Feb2020 records, 10-fold CV: ACC
90.33, MCC 0.8076, AUC 96.36 versus published 89.9/0.799/96.2 ---
descriptive parity with a slight edge, no significance claim possible.
\item[Study 15 (ablation).] Removing the P0ACW4 signal peptide drops the
ensemble score from 0.9082 to 0.6299: a negative control showing the
score partly tracks trafficking signal.
\item[Study 16 (homology).] 15.98\% of held-out Feb2020 sequences share
an exact decamer with training; random-CV results are interpolation.
\item[Study 17 (strata).] The large-dataset test AUROC is 0.9206 overall,
0.9643 on the tiny reviewed-positive subgroup, 0.9121 on the confounded
majority: the headline number cannot be attributed to AMP function.
\item[Study 18 (composition).] First-order composition barely separates
the label $\times$ review-status groups; the learnable confound is mostly
a length artifact, directing the fix toward length-matched reviewed-only
redesign.
\end{description}
